\chapter{Conclusion}

\section{Thesis Summary}

In this thesis I have used direct numerical simulations to study the attractor states of small-scale dynamos (SSDs), and quantified how compressibility shifts attractor properties. While the results in \S\ref{chapter:ssd-exp} and \S\ref{chapter:ssd-nl} have broader implications for SSDs in $\mathrm{Pm} \geq 1$ plasmas, I use galactic magnetism as the lens for this summary, and refer the reader back to the concluding sections of each paper-chapter for a wider-scope discussion and summary of results.

In \S\ref{chapter:ssd-exp}, I explored how large-scale magnetic structure could be generated in young galaxies, long before large-scale coherence is usually expected to develop. Using periodic-box simulations of controlled and isolated SSD amplification, I quantified how compressive turbulence modifies the magnetic structures produced by SSD action. A key challenge of this process was that several measurements of field-line morphology are sensitive to dissipation and therefore to numerical resolution. To overcome this, I developed a statistical framework to extract attractor properties that give a robust, resolution-independent picture of the kinematic phase. Using this framework, I ran the most comprehensive suite of $ \mathrm{Pm} \geq 1 $ SSD simulations to date, spanning a wide range of plasma conditions, from viscous and smooth subsonic flows through to turbulent, highly supersonic, shock-dominated flow regimes. This allowed me to isolate how shocks and compressive modes modify the characteristic magnetic length scales. For galactic magnetism, the key takeaway is that while compressibility can shift the magnetic correlation scale to much larger values than incompressible SSD theory predicts, it does not by itself explain the $\sim \mathrm{kpc}$ coherence inferred in young galaxies. Instead, compressive turbulence provides a natural pathway to generate $\sim \mathrm{pc}$ fields, which must be further organised to system scales by other system-level processes.

In \S\ref{chapter:ssd-nl}, I turned to the more dynamically relevant, nonlinear stage of SSD evolution. In this regime I showed that magnetic amplification, which is intrinsically transient\footnote{
    This phase mediates the transition from the kinematic SSD attractor, where the flow is effectively quasi-hydrodynamic (only kinetic stresses affect the dynamics), to a maintained MHD regime in which the dynamics are regulated by both magnetic and kinetic stresses.
}, becomes liberated from dissipation regulated effects. In highly compressible turbulence this transient phase is difficult to identify because stochastic fluctuations obfuscate the underlying trend. To separate the underlying evolution from this variability, I developed a hierarchical Bayesian framework that couples many repeated experiments of the same setup. With this approach, I showed that the fraction of the turbulent energy injection converted into magnetic energy is approximately constant across all turbulent regimes, and is consistent with the efficiency during the kinematic phase. An immediate consequence of this is an as-yet unexplained universality in the duration of nonlinear growth when measured in turbulent eddy turnover times. In the context of galaxy evolution, this implies that even the nonlinear SSD can re-establish a maintained magneto-hydrodynamic (MHD) state on a relatively short dynamical timescale following merger- or fly-by-driven disruption.

In both \S\ref{chapter:ssd-exp} and \S\ref{chapter:ssd-nl}, I have deliberately explored SSD physics in isolation. This was necessary to concentrate numerical resolution on the inertial and resistive scales that control SSD amplification, and to develop statistical tools that extract robust measurements from the stochastic processes at play. Galactic magnetism, however, ultimately requires system-scale dynamics. If compressive turbulence in young, rapidly evolving galaxies provides a natural route to generating a larger-scale seed field early, then the remaining step is to identify which global processes subsequently act on that seed to build the $\sim \mathrm{kpc}$ coherence we see. Answering this requires simulations that resolve \textit{both} the SSD-scale physics and \textit{also} the larger-scale environment that organises the field. This is challenging with current CPU-based codes, but becomes increasingly feasible on modern GPU-accelerated supercomputers, especially using codes like \textsc{Quokka}, which were designed architecturally from the outset for these hardware capabilities. To this end, in \S\ref{chapter:quokka} I contributed to the development of the new radiation-MHD code \textsc{Quokka}, and implemented a new constrained-transport MHD module that maintains $\nabla\cdot\mVector{b} = 0$ to numerical precision by construction. I both implemented a number of existing state-of-the-art schemes and introduced a new approach in the paper, which proved to be the best scheme for the performance requirements of \textsc{Quokka}. The MHD module in \textsc{Quokka} now enables me and others to explore the next questions about galactic magnetism, which will require resolution sufficient to capture SSD amplification while simultaneously resolving the system-scale mechanisms that can organise the resulting seed field. In the following section I will highlight some of these questions.

\section{Future Work}

The MHD module that I implemented in \textsc{Quokka} has already been put to work in two projects I am involved with. One was recently awarded over a million node hours on \textsc{Frontier} and will run one of the highest-resolution isolated galaxy simulations to-date, with the goal of exploring the phase structure and properties of galactic winds. The other project focuses on determining what sets the magnetic field reversal scale in Milky Way-like galaxies. Moreover, building on the work in this thesis, I plan to use \textsc{Quokka} both to pursue scientific questions and to extend its MHD capabilities. Below I briefly outline these directions and the code developments required to enable them.

\subsection{The Saturation of the Small Scale Dynamo}

The simulations in \S\ref{chapter:ssd-exp} and \S\ref{chapter:ssd-nl} established robust attractor properties of the kinematic and nonlinear SSD, including resolution-independent morphology and a universal nonlinear growth efficiency (through the lens of turbulent eddy turnover times). Across the parameter space explored, I did not find evidence that magnetic dissipation altered the qualitative flow regime; however, my simulations were confined to modest $\mathrm{Re}$ and $\mathrm{Rm}$, where dissipation remained dynamically important and the kinematic phase was explicitly Ohmic-limited. Most astrophysical plasmas are in the $\mathrm{Rm} \gg 1$ regime, though, where dissipation is expected to become confined to increasingly thin current sheets, with magnetic energy redistributed through intermittent structures that can become unstable to tearing \citep{loureiro2017role, galishnikova2022tearing, beattie2025spectrum}. It therefore remains unclear whether the SSD approaches a dissipation-independent attractor in this $\mathrm{Rm}$ limit, or whether the maintained MHD state continues to depend on microphysical scales. To address this, I will extend \textsc{Quokka} beyond the MHD regime by implementing explicit (shear and bulk) viscosity to sit alongside Ohmic resistivity, enabling controlled variation of $\mathrm{Pm}$ and allowing the physical dissipation scale to be set and systematically pushed to smaller values while remaining resolved (as was done in this thesis). Combined with the GPU-native performance of \textsc{Quokka}, this will make substantially higher effective $\mathrm{Re}$ and $\mathrm{Rm}$ attainable than were accessible in this thesis \citep{ji2022magnetic}, and will allow a direct test of whether tearing-mediated dynamics begin to regulate saturation as current sheets become sufficiently thin. The goal is to determine how the magnetic and kinetic energy spectra, and the maintained field geometry, evolve through this transition, and whether the nonlinear SSD converges toward an asymptotic configuration.

\subsection{Magnetic Field Line Curvature}

High-energy cosmic rays (CRs) are observed to isotropise efficiently, and this has been hypothesised to be a consequence of the geometry of the magnetic field lines that thread the ISM. In these environments, CRs stream along field lines and scatter when they encounter curvature events on scales comparable to their gyroradius \citep{kempski2023cosmic, Lemoine2023_field_reversal_CR_transport}. A simple criterion for this process relies on the empirical relation $\kappa \propto b^{1/2}$ between magnetic field strength $b \equiv (\mVector{b}\cdot\mVector{b})^{1/2}$ and line curvature $\kappa = |(\mVectorUnit{b}\cdot\nabla)\mVectorUnit{b}|$. In \S\ref{chapter:ssd-exp}, I showed that this scaling emerges naturally during the kinematic phase of the SSD, consistent with earlier theoretical arguments by \citet{schekochihin2004simulations}, which showed that the quantity $b \kappa^{\alpha}$ evolves according to
\begin{align}
    \mDDFrac{\ln(b \kappa^\alpha)}{t}
        = \frac{1}{b}\mDDFrac{b}{t} + \frac{\alpha}{\kappa}\mDDFrac{\kappa}{t}
        = \rbrac{
                \alpha (\mVectorUnit{\kappa}\otimes\mVectorUnit{\kappa})
                + (1 - 2\alpha) (\mVectorUnit{b}\otimes\mVectorUnit{b})
                - \mTensor{I}
            }:\nabla\otimes\mVector{u},
\end{align}
with turbulent plasmas organising toward $\alpha = 1/2$. This behaviour has been seen across a wide range of magnetised systems: from incompressible SSDs \citep{schekochihin2004simulations}, mildly compressible maintained MHD states \citep{yuen2020curvature}, the Solar wind \citep{huang2020observations, bandyopadhyay2020situ, hu2025interplanetary}, and the Earth's magnetosheath \citep{ji2022curvature}. The ubiquity of this result suggests that the curvature distribution reflects a universal geometric constraint of magnetised flows. However, its origin remains unclear, and it is not known how robust the $\kappa \propto b^{1/2}$ coupling is once compressibility is important. In \S\ref{chapter:ssd-exp}, I found that increasing compressibility weakens this dependence, indicating that shocks and compressive modes modify the geometric relationship between field strength and curvature even during the kinematic phase. Understanding what sets this deviation requires a more detailed investigation of the curvature statistics and their connection to the underlying flow structures. I therefore plan to quantify the curvature distribution and its coupling to field strength in the nonlinear and high-$\mathrm{Rm}$ regimes accessible with \textsc{Quokka}, and to assess how these changes in magnetic geometry impact CR scattering and transport.

\subsection{The Journey From Small to Large Scale Dynamos During Galaxy Mergers}

Galaxy mergers are a setting in which SSDs and large-scale dynamos can operate in sequence. Merger-driven turbulence rapidly amplifies magnetic energy through SSD action, while the post-merger remnant introduces rotation, shear, and stratification that can organise fluctuating fields into large-scale coherence. Existing numerical studies already explore aspects of this scenario, but they largely capture different parts of the problem. Tall-box simulations \citep[\eg][]{gent2023transition} follow mean-field growth in controlled domains and show that stochastic fluctuations can slow the emergence of coherence, but their restricted geometry limits the large-scale modes that can grow, and necessarily ignore topological constraints that can only be captured in simulations including full galactic discs. Conversely, full galaxy merger simulations \citep[\eg][]{whittingham2021impact, whittingham2023impact} resolve efficient SSD amplification in the dense, turbulent central remnant, and show that differential rotation spreads magnetic energy across the disc post-merger, but their resolution is insufficient to follow the subsequent organisation into a true large-scale dynamo. My goal is to bridge these approaches using \textsc{Quokka} in an intermediate setup: two isolated galaxies (free of cluster-scale forcing) simulated at sufficient resolution to capture SSD amplification while retaining the full galactic dynamics required for large-scale organisation. This class of setup has been used extensively to study galaxy formation, but, to the best of my knowledge, has not been used to quantify dynamo action across the transition from SSD to LSD dominated dynamics. With GPU-accelerated performance, this becomes a tractable route to measure how quickly, and by which mechanisms, merger-driven SSD fields transition into coherent, system-scale structure in the post-merger remnant.
